% Note to reader: lines beginning with the '%' character are
% 'comments' to you, the human reader of this code, and are
% ignored by the LaTeX compiler.

%%%%%%%%%%%%%%%%%%%%%%%%%%%%%%%%%%%%%%%%%%%%%%%%%%%%%%%%%%%%%%%%%%
% = Sample template for MIT Junior Lab Student Written Summaries =
% Available from
%    http://web.mit.edu/8.13/www/Samplepaper/sample-paper.tex
%
% Last Updated July 23, 2014
%
% Adapted from the American Physical Societies REVTeK-4.1 Pages
%    at http://publish.aps.org
%
% ADVICE TO STUDENTS: Each time you write a paper, start with this
%    template and save under a new filename.  If convenient, don't
%    erase unneeded lines, just comment them out using the
%    '%' character at the start of the line.  Often, they
%    will be useful containers for information.
%
% Using pdflatex, images must be either PNG, GIF, JPEG or PDF.
%    Turn EPS (encapsulated postscript) images to PDF using the
%    epstopdf utility on UNIX systems.
%%%%%%%%%%%%%%%%%%%%%%%%%%%%%%%%%%%%%%%%%%%%%%%%%%%%%%%%%%%%%%%%%%

%%%%%%%%%%%%%%%%%%%%%%%%%%%%%%%%%%%%%%%%%%%%%%%%%%%%%%%%%%%%%%%%%%
% = TO COMPILE THIS DOCUMENT =
%
% From the command line, it would go like this --- assuming you are
%    in the directory where the filename.tex source file and the
%    filename.bib bibliography file are located, and that you have
%    permission to create and write files in that directory:
%      > pdflatex filename
%      > bibtex filename
%      > pdflatex filname
%      > pdflatex filename
%    Yes, you run the command several times. The earlier runs
%    create auxilliary files which keep track of references,
%    citations, equation and section numberring, etc. The later
%    runs combine the information in these auxilliary files with
%    your source document to create the finished product.
%
% If you are using a GUI LaTeX editor like TeXWorks, then there
%    is probably a menu bar button for pdfLaTeX and another for
%    BibTeX. Hit them in the order indicated above. There is
%    probably also a 'TeXify' button, or something similarly named,
%    which runs all the above commands in one shot.
%%%%%%%%%%%%%%%%%%%%%%%%%%%%%%%%%%%%%%%%%%%%%%%%%%%%%%%%%%%%%%%%%%


%%%%%%%%%%%%%%%%%%%%%%%%%%%%%%%%%%%%%%%%%%%%%%%%%%%%%%%%%%%%%%%%%%
%  = PREAMBLE =
% The preamble of a LaTeX document is the set of commands that precede
% the \begin{document} line.  It contains a \documentclass line
% to load the REVTeK-4.1 macro definitions and various \usepackage
% lines to load other macro packages.
%
% ADVICE TO STUDENTS: This preamble contains a suggested set of
%     class options to generate a ``Junior Lab'' look and feel that
%     facilitate quick review and feedback from one's peers, TAs,
%     and section instructors.  Don't make substantial changes
%     to the style without first consulting your section
%     instructor.
%%%%%%%%%%%%%%%%%%%%%%%%%%%%%%%%%%%%%%%%%%%%%%%%%%%%%%%%%%%%%%%%%%

%\documentclass[aps,twocolumn,secnumarabic,balancelastpage,amsmath,amssymb,nofootinbib, floatfix]{revtex4}
\documentclass[aps,twocolumn,secnumarabic,balancelastpage,amsmath,amssymb,nofootinbib,floatfix]{revtex4-1}

%%%%%%%%%%%%%%%%%%%%%%%%%%%%%%%%%%%%%%%%%%%%%%%%%%%%%%%%%%%%%%%%%%%
% N.B.:  Different computers have different packages installed.
%        To compile this template in the current Athena
%        environment, REVTeX 4.1 must be used.  To use the older
%        REVTeX 4, switch which documentclass line above is
%        commented out above. There are ``bad'' distributions of
%        LaTeX for Windows available on the internet which may
%        cause users to struggle unjustifiably with REVTeX 4.1.
%
%        If you are unable to compile the template at all, you
%        may need to update your LaTeX packages. (Alternatively, if
%        your LaTeX distribution includes only the older RevTEX 4,
%        then try changing the documentclass line above. In particular,
%        this approach solves a common compilation problem for users of
%        the TeXWorks editor on Windows, which presents erroneously as a
%        error in the bibliography file.) Don't hesitate to speak
%        with your section instructor or a TA if you're having
%        issues getting this template to compile.
%%%%%%%%%%%%%%%%%%%%%%%%%%%%%%%%%%%%%%%%%%%%%%%%%%%%%%%%%%%%%%%%%%%

%%%%%%%%%%%%%%%%%%%%%%%%%%%%%%%%%%%%%%%%%%%%%%%%%%%%%%%%%%%%%%%%%%%
% = Explanation of documentclass options =
%
% aps, prl stand for American Physical Society and Physical
%     Review Letters respectively.
% twocolumn permits two columns, of course.
% nobalancelastpage doesn't attempt to equalize the lengths of
%     the two columns on the last page  as might be desired in a
%     journal where articles follow one another closely.
% amsmath and amssymb are necessary for the subequations
%     environment among others. These functionalities can
%     also be added use the usepackage function described below,
%     but REVTeX conveniently includes them as documentclass
%     options.
% secnumarabic identifies sections by number to aid electronic
%     review and commentary.
% nofootinbib forces footnotes to occur on the page where they are
%      first referenced and not in the bibliography.
% floatfix attempts to help LaTeX decide where to place ``floats'',
%      like figures and plots, when it gets stuck and can't decide
%      by it's normal algorithm.
% REVTeX 4.1 is a set of macro packages designed to be used with
%      LaTeX 2e. REVTeX is well-suited for preparing manuscripts
%      for submission to APS journals.
%
% = Other documentclasses =
%
% The 'revtex4' and 'revtex4-1' documentclasses are somewhat
%    specialized for making documents in the style of the APS
%    journals. For a more standard or generic looking LaTeX paper,
%    you could try any of the built-in documentclasses, in
%    particular 'article' or 'report'. Someday, you may try to use
%    the 'mitthesis'  documentclass available for download from the
%    MIT Libraries. The vast majority of source code written for
%    one documentclass should work just fine in any other, but
%    occasional quirks arise. For example, some documentclasses
%    disagree on whether the abstract declaration should come
%    before or after the \begin{document} declaration.
%
%%%%%%%%%%%%%%%%%%%%%%%%%%%%%%%%%%%%%%%%%%%%%%%%%%%%%%%%%%%%%%%%%%%

%% Now, include some packages which provide new commands that
%% extend LaTeX's capabilities. Note that the nearly-essential
%% AMS math packages were added already as documentclass options
%% for REVTeX, but could have been added here using
%% \usepackage{amsmath}, etc. The pacakges below are commonly
%% useful, but there are many, many more available to solve a
%% multitude of typesetting quandries (google your problem),
%% and you  probably have the necesary packages installed on your
%% system already. Among the examples listed below, this sample
%% document only actually makes use of the 'graphicx', 'bm',
%% and 'hyperref' pacakges, so the others are commented out for
%% tidyness.


\usepackage{graphicx}      % tools for importing graphics
%\usepackage{lgrind}        % convert program code listings to a form
                            % includable in a LaTeX document
%\usepackage{xcolor}        % produces boxes or entire pages with
                            % colored backgrounds
%\usepackage{longtable}     % helps with long table options
%\usepackage{epsf}          % old package handles encapsulated postscript issues
\usepackage{bm}            % special bold-math package. usge: \bm{mathsymbol}
%\usepackage{asymptote}     % For typesetting of mathematical illustrations
%\usepackage{thumbpdf}
\usepackage[colorlinks=true]{hyperref}  % this package should be added after
                                        % all others.
                                        % usage: \url{http://web.mit.edu/8.13}


%%%%%%%%%%%%%%%%%%%%%%%%%%%%%%%%%%%%%%%%%%%%%%%%%%%%%%%%%%%%%%%%%%%
% And now, begin the document...
%%%%%%%%%%%%%%%%%%%%%%%%%%%%%%%%%%%%%%%%%%%%%%%%%%%%%%%%%%%%%%%%%%%

\begin{document}
\title{Title}
\author{Raul Hernandez Sosa}
\email{raulh482@mit.edu}
%\homepage{http://web.mit.edu/8.13/} %If you don't have one, just comment out this line.
\date{\today}
\affiliation{MIT Department of Physics}


\begin{abstract}
    We investigated the momentum–velocity relationship of relativistic electrons, enabling a comparison between relativistic and nonrelativistic dynamics along with a determination of the electron charge-to-mass ratio.
    Electrons from a strontium-90/yttrium-90 source were subjected to magnetic and electric fields within an apparatus of known geometry, allowing for measurement of electron momentum and velocity.
    The relativistic model provided a better fit to the data and a charge-to-mass ratio of $(1.515\pm0.004\,\mathrm{stat}\pm0.127\,\mathrm{syst})\times10^{11}\,\mathrm{C\,kg^{-1}}$.
    Statistical and systematic uncertainties were evaluated to assess the significance of the discrepancy between the measured and reference electron charge-to-mass ratios.
\end{abstract}

\maketitle
\raggedbottom





\section{Introduction}
By the early twentieth century, a conflict had emerged between Maxwell's equations, the principle of relativity, and
the classical treatment of velocity. Maxwell's equations predict the existence of electromagnetic waves that propagate at speed $c$.
The principle of relativity posits that physical laws are invariant between inertial reference frames. Together, these two assumptions predict that
electromagnetic waves should propagate at $c$ in every inertial reference frame. In contrast, classical velocity addition implies that
the measured speed of an electromagnetic wave by an observer should depend on that observer's inertial reference frame.

A resolution to this conflict came in the form of what is now known as Einstein's theory of special relativity~\cite{einstein1905}.
%This theory posits that
%to properly convert between inertial reference frames, separate notions of space and time should be abandoned in favor of a combined space-time.
Special relativity predicts a relativistic momentum:
$$\bm{p}=\gamma m\bm{v}$$
where
$\gamma =1/\sqrt{1- v^2/c^2}$. The relativistic momentum approaches the nonrelativistic momentum:
$$\bm{p}=m\bm{v}$$
if $v\ll c$.

At velocities approaching $c$, the momentum-velocity relationship can be used to experimentally test special relativity.
The purpose of this experiment is to observe the momentum-velocity relationship of electrons moving at relativistic speeds,
and to compare this observation to the predictions of relativistic and nonrelativistic dynamics. To do this, we make use of electrons
subject to the Lorentz force. This approach allows for a determination of the electron charge-to-mass ratio for both the relativistic
and nonrelativistic regimes.




\section{Experiment}
A schematic of the experimental apparatus for measuring relativistic electrons
can be seen in Figure~\ref{fig:apparatus}.
The apparatus is made up of a vacuum chamber surrounded by
an electromagnet.

\begin{figure}[htbp]
    \centering
    \includegraphics[
        page=3,
        trim={20mm 210mm 110mm 15mm},
        clip,
        width=\columnwidth
    ]{JLExp09.pdf}
    \caption{Schematic of the electron trajectory and velocity selector, reproduced from the laboratory manual~\cite{labmanual}.}
    \label{fig:apparatus}
\end{figure}


Within the vacuum chamber, there is a sample of the radioactive isotope strontium-90,
which passively decays into yttrium-90. In turn, yttrium-90 decays into stable
zirconium-90. Both of these decay processes produce electrons with a distribution of velocities.


Also within the vacuum chamber, there is a 'velocity selector' placed in the same plane as the electron source and at a distance of $2\rho$ away from it.
The velocity selector is made up of two conducting plates with a variable voltage $V$ applied between them.

The electromagnet is made up of a spherical arrangement of currents such that
an approximately uniform magnetic field, $\bm{B}=B\hat{z}$, is created within its volume.
Electrons produced at the source are subject to the Lorentz force:
$$\bm{F}=-e(\bm{E}+\bm{v} \times \bm{B})$$
where $\bm{B}=B\hat{z}$ and $\bm{v}$ is the velocity of a produced electron.

Outside of the velocity selector, $\bm{E}=0$, thus in this region, electrons are subject to a net force with magnitude
$|\bm{F}|=evB$ in the direction perpendicular to that of their velocities. Therefore, electrons produced at the source
adopt cyclotron orbits, with orbital radii dependent on the speed at which they
were produced from their originating decay event.

A barrier is also present within the vacuum chamber. This barrier has a slit such that only electrons with orbital
radii close to $\rho$ are able to reach the velocity selector.

Electrons in cyclotron orbits have momenta that are constant in magnitude, but variable in direction
as per their angular speed $\omega$. Thus, electrons that make it to the velocity selector can be written to have:
$\left\lvert d\bm{p}/dt\right\rvert=pv/\rho$.
By relating the net force acting on a produced electron to its change in momentum while it is outside the velocity selector, we find:
$$p=e \rho B$$
Therefore, with a constrained orbital radius $\rho$, the choice of $B$ selects the momentum of the electrons that reach
the velocity selector. Comparing this result for momentum to the nonrelativistic and relativistic forms of momentum, we can find:
$v=(e/m)B\rho$ and $v\gamma =(e/m) B\rho$ respectively.


The velocity measurement of electrons in the apparatus involves the velocity selector. Electrons moving through
the velocity selector are subject to a Lorentz force: $|\bm{F}|=e|E-vB|$ perpendicular to their velocities.
Sufficiently deflected electrons strike the plates, while electrons satisfying $E=vB$ ideally pass through without deflection.

The electric field between the plates of the velocity selector is related to the
potential difference $V$ between the plates by: $|\bm{E}|=V/d$, where $d$ is the plate spacing. Thus, for electrons that
make it through the velocity selector we can find:
$$v=\frac{V}{dB}$$

After the velocity selector, there is a PIN diode, which records the count rate of electrons striking its surface. The diode has a semiconductor layer in which incident electrons
create electron--hole pairs, producing electrical pulses that
are counted.



The finite length of the velocity selector makes it so that electrons can reach the PIN diode when $V \neq vBd$. To find the relevant
$V_0$ such that $v=V_0/(dB)$, we measure electron count rate at a range of voltages at a set magnetic field, and determine $V_0$ to be the
voltage at which we observe the maximal count rate (more on the exact procedure for our determination of $V_0$ is provided in the next section).



The procedure used to collect $p$ and $v$ measurements was as follows:
\begin{enumerate}
\item  Place a calibrated Hall effect magnetometer at a known position ('position 1') within the electromagnet directly above the electron trajectory.
\item  Adjust the voltage and current of electromagnet to reach a desired magnetic field.
\item  Keeping magnetic field fixed, vary voltage of velocity selector by $\pm1\,\mathrm{kV}$ centered at the expected
$V_0$ (given by relativistic dynamics and the parameters of the apparatus). Record total electron counts at each voltage for 240 seconds. Collect
at least 5 data points of counts vs. voltage per magnetic field.
\item Repeat (1.)-(3.) for a total of eight magnetic fields ranging from 70 to 115 gauss.
\end{enumerate}
During our data collection, magnetic fields of 70, 80, 90, 95, 100, 105, 110, and 115 gauss were used.


\section{Results and Analysis}
Figure~\ref{fig:counts} is a plot of counts as a function of velocity selector voltage for a magnetic field of 70 gauss, collected as described in the above procedure.
Electron counting was treated as a Poisson process, giving a statistical uncertainty of $\sqrt{N}$ for a total count $N$. To determine $V_0$, the count-versus-voltage
data were fitted with a Gaussian plus a constant background using $\chi^2$ minimization.
\begin{figure}[htbp]
    \centering
    \includegraphics[
        width=\columnwidth
    ]{Figure 1.pdf}
    \caption{Total counts versus selector voltage at 70 gauss. The curve is a Gaussian plus constant-background fit. Error bars represent Poisson counting uncertainties.}
    \label{fig:counts}
\end{figure}

The equilibrium voltage $V_0$ is given by the fitted Gaussian peak position,
and its statistical uncertainty is obtained from SciPy's \texttt{curve\_fit} fitting procedure.
Using this process, we determine $V_0$ for each magnetic field, which, along with the parameters of the apparatus, allows us to calculate the electron velocity.



If the nonrelativistic regime properly describes the momentum-velocity relationship of electrons in this apparatus, we expect
$\beta=\frac{e}{mc}B\rho$, where $\beta=v/c=V_0/(dBc)$. Thus, we expect a linear relationship between $\beta$ and $B$ in the nonrelativistic regime. Figure~\ref{fig:beta_vs_B} shows
$\beta$ as a function of $B$. We fit a line of the expected form $\beta=kB$ through the origin using $\chi^2$ minimization, weighted by the statistical
uncertainties in $\beta$. The fitted slope $k$ gives the nonrelativistic charge-to-mass ratio through $e/m=ck/\rho$.

\begin{figure}[htbp]
    \centering
    \includegraphics[
        width=\columnwidth
    ]{Figure 5.pdf}
    \caption{$\beta$ versus magnetic field. The solid line is the nonrelativistic fit through the origin. Dotted curves show the nonrelativistic and relativistic predictions using the reference value of $e/m$. Error bars are statistical only.}
    \label{fig:beta_vs_B}
\end{figure}

If the relativistic regime properly describes the momentum-velocity relationship of electrons in this apparatus, we expect
$\beta\gamma=\frac{e}{mc}B\rho$, where $\beta=v/c=V_0/(dBc)$. Thus, we expect a linear relationship between $\beta\gamma$ and $B$ in the relativistic regime. Figure~\ref{fig:beta_gamma_vs_B} shows
$\beta\gamma$ as a function of $B$. We fit a line of the expected form $\beta\gamma=kB$ through the origin using $\chi^2$ minimization, weighted by the statistical
uncertainties in $\beta\gamma$. The fitted slope $k$ gives the relativistic charge-to-mass ratio through $e/m=ck/\rho$.

\begin{figure}[htbp]
    \centering
    \includegraphics[
        width=\columnwidth
    ]{Figure 4.pdf}
    \caption{$\beta\gamma$ versus magnetic field. The solid line is the relativistic fit through the origin. Dotted curves show the predictions of both models using the reference value of $e/m$. Error bars are statistical only.}
    \label{fig:beta_gamma_vs_B}
\end{figure}

The data fitted to the nonrelativistic regime produced a reduced $\chi^2$ of 246.01, while the relativistic regime produced a reduced $\chi^2$ of 31.10. This indicates that the relativistic regime provides a better description of the data.

These fitting procedures provide values of $e/m$ for both the relativistic and nonrelativistic models, as summarized in Table~\ref{tab:em_final} alongside the reference value \cite{nist_em}.
The statistical uncertainty in $V_0$ was propagated into $\beta$ and $\beta\gamma$ for the nonrelativistic and relativistic models, respectively.
These uncertainties were used to weight the linear fits and determine the uncertainty in each fitted slope. The slope uncertainty was then propagated
through $e/m=ck/\rho$ to obtain the statistical uncertainty in $e/m$ for each regime. Systematic uncertainties are evaluated separately.

As a consistency check, we calculated $e/m$ separately at each magnetic
field using $e/m=v/(B\rho)$ for the nonrelativistic model and
$e/m=\gamma v/(B\rho)$ for the relativistic model.
Figure~\ref{fig:em_vs_B} shows these estimates with statistical error
bars. Our reported overall values of $e/m$ remain those obtained
from the fitted slopes.

\begin{figure}[htbp]
    \centering
    \includegraphics[width=\columnwidth]{Figure 6.pdf}
     \caption{Individual estimates of $e/m$ versus magnetic field
    for the relativistic and nonrelativistic models. Error bars
    represent statistical uncertainties propagated from $V_0$.
    The horizontal dotted line indicates the reference value.}
    \label{fig:em_vs_B}
\end{figure}




The nonrelativistic estimates of $e/m$ generally decrease with increasing
magnetic field, whereas the relativistic estimates show less field
dependence. Since the electron charge-to-mass ratio should be independent
of velocity, this behavior favors the relativistic description.












We evaluated three sources of systematic uncertainty: magnetic field nonuniformity, the $\rho$ measurement and the $d$ measurement.
To quantify the systematic uncertainty in our magnetic field measurement due to our assumption of a uniform magnetic field, we measured the magnetic field at 5 known positions (positions 1-5) along the
electron trajectory. The exact procedure was as follows:
\begin{enumerate}
\item  Place a calibrated Hall effect magnetometer at position 1. Adjust voltage and current of the electromagnet to produce a desired magnetic field. Record the values of voltage and current.
\item  Move the magnetometer to another of the 5 known positions. Adjust voltage and current of the electromagnet to the settings recorded in (1.). Record the magnetic field.
\item  Repeat (2.) until magnetic field is measured at all 5 known positions.
\item  Repeat (1.)-(3.) for all 8 magnetic fields used in the $p$, $v$ measurements.
\end{enumerate}


\begin{figure}[htbp]
    \centering
    \includegraphics[
        width=\columnwidth
    ]{Figure 2.pdf}
    \caption{Magnetic field versus probe position, with the field at position 1 set to 95 gauss.}
    \label{fig:B_vs_position}
\end{figure}


Figure~\ref{fig:B_vs_position} shows the measured magnetic field as a function of position along the electron trajectory when the magnetic field at position 1 is 95 gauss.

At each setting, we calculated the sample standard deviation of the five
field measurements and divided it by their mean to obtain a fractional
variation. Averaging the eight fractional variations with equal weights
gave $2.3\%$. We adopted this as a common
fractional systematic uncertainty scale for the magnetic field.

The geometric uncertainties were taken from the orbit radius,
$\rho=(0.203\pm0.002)\,\mathrm{m}$, and plate separation,
$d=(1.80\pm0.03)\,\mathrm{mm}$~\cite{labmanual}.

These input uncertainties were propagated into the charge-to-mass ratios
obtained from the fitted slopes. For each source, the input was shifted
upward and downward by its uncertainty, the derived quantities and fit
weights were recalculated, and the linear fit was repeated. The larger
absolute change in $e/m$ was retained as that source's contribution.
The magnetic-field shifts were applied coherently to all eight settings,
treating the adopted fractional field uncertainty as a shared
systematic scale.

The three contributions were combined in quadrature, assuming independent
sources. Table~\ref{tab:systematic_budget} lists the resulting fractional
uncertainties in $e/m$. The different contributions for the two models
arise from their different dependence on the measured quantities,
including the additional velocity dependence of $\gamma$ in the
relativistic model.


\begin{table}[htbp]
    \centering
    \begin{tabular}{lcc}
        \hline
        Source & Relativistic & Nonrelativistic \\
        \hline
        Magnetic field   & $7.49\%$ & $4.69\%$ \\
        Plate separation & $3.65\%$ & $1.69\%$ \\
        Orbit radius     & $1.00\%$ & $1.00\%$ \\
        \hline
        Total systematic & $8.39\%$ & $5.09\%$ \\
        \hline
    \end{tabular}
    \caption{Systematic uncertainty contributions to $e/m$}
    \label{tab:systematic_budget}
\end{table}

Including these systematic estimates, the fitted-slope measurements
are summarized in Table~\ref{tab:em_final}.

\begin{table}[htbp]
    \centering
    \begin{tabular}{lccc}
        \hline
        Model & $e/m$ & Statistical & Systematic \\
        \hline
        Relativistic    & $1.515$ & $\pm0.004$  & $\pm0.127$ \\
        Nonrelativistic & $1.034$ & $\pm0.0013$ & $\pm0.053$ \\
        Reference      & $1.75882$ & --- & --- \\
        \hline
    \end{tabular}
    \caption{Charge-to-mass ratios obtained from the fitted slopes,
    with statistical and systematic uncertainties.
    All numerical entries are in units of
    $10^{11}\,\mathrm{C\,kg^{-1}}$.}
    \label{tab:em_final}
\end{table}

An independently conducted analysis by my laboratory partner determined
$V_0$ using counts selected from an energy-spectrum peak rather than
the full recorded spectrum. As shown in Figure~\ref{fig:partner_comparison},
this analysis yielded similar velocities and also favored the relativistic
model, supporting our conclusion. Both analyses use the same underlying
measurements and therefore share experimental systematic uncertainties.

\begin{figure}[htbp]
    \centering
    \includegraphics[width=\columnwidth]{Partner_comparison.pdf}
    \caption{Figure reproduced from my laboratory
    partner's analysis. The axis notation $\beta=E/B$ corresponds to Gaussian
    cgs units; in SI units, $\beta=E/(Bc)$.}
    \label{fig:partner_comparison}
\end{figure}


\section{Conclusion}

We measured the momentum--velocity relationship of electrons selected over
magnetic fields from 70 to 115 gauss. Magnetic deflection determined the
selected momentum, while the peak transmission voltage of a crossed-field
selector determined the velocity.

The relativistic model described the measured momentum--velocity
relationship better than the nonrelativistic model, producing a
lower reduced $\chi^2$ and less field dependence in the individual
$e/m$ estimates. The fitted relativistic result,
$(1.515\pm0.004\,\mathrm{stat}\pm0.127\,\mathrm{syst})
\times10^{11}\,\mathrm{C\,kg^{-1}}$,
remained below the reference value by approximately 1.9 times
the estimated systematic uncertainty. Finite trajectory acceptance
may contribute additional systematic shifts because electrons
follow different radii and sample different magnetic fields.
The results favor relativistic dynamics, although the remaining
discrepancy limits the precision of this test.

\section*{Acknowledgments}
AI statement: AI was used to develop plotting scripts, and for questions about the uncertainty process. Codes and ideas were human reviewed.

\begin{thebibliography}{9}

\bibitem{labmanual}
MIT Department of Physics,
\textit{Relativistic Dynamics},
Junior Laboratory Experiment 09.

\bibitem{nist_em}
National Institute of Standards and Technology,
\textit{CODATA Value: Electron Charge to Mass Quotient},
2022 CODATA adjustment,
\url{https://physics.nist.gov/cgi-bin/cuu/Value?esme}.

\bibitem{einstein1905}
A. Einstein, Ann. Phys. \textbf{17} (1905).
\end{thebibliography}





\end{document}
